\section{The complete Liouville count for arbitrary phases}
\label{sec:integrability-proof}

We prove \cref{thm:canonical-independence} using the commuting algebra
$\mathcal A$ defined in \cref{sec:canonical-independence}.
The argument inserts one index at a time; spectral invariants and
equal-phase moment invariants together supply the new independent
functions.

The next lemma constructs an independence certificate in the mode
coordinates. A vector $v$ is \emph{cyclic} for a matrix $\mathsf Z$ of size
$a$ if $v,\mathsf Z v,\ldots,\mathsf Z^{a-1}v$ span $\C^a$.

\begin{lemma}[Increasing a matrix dimension]\label{lem:canonical-insertion}
Increase one of $L,M,N$ by one. Suppose the new phase already
occurs $r$ times in that group and $c$ other indices there have
different phases. Let the other group sizes be $m_1,m_2$, and put
$n=m_1+m_2$. Assume either $r\geq n$ or
\begin{equation}\label{eq:canonical-insertion-condition}
 |m_1-m_2|\leq r+c+1.
\end{equation}
If the old algebra $\mathcal A$ supplies the complete Liouville count,
the enlarged algebra does too. The insertion supplies
$n-\min(r,n)$ new independent spectral functions and $\min(r,n)$
new independent functions of the principal moments.
\end{lemma}
\begin{proof}
Write the new outgoing entries as a row $\boldsymbol a\in\C^{1\times n}$ and
the incoming entries as a column $\boldsymbol b\in\C^{n\times1}$.
We test the Jacobian at $\boldsymbol a=0$. The enlarged determinant then equals
$(\lambda-e-s\nu)F(\lambda,\nu)$, where $e,s$ are the new intercept
and slope. More generally the bordered determinant is this product
minus a term bilinear in the new outgoing and incoming entries.
Multiplication by $\lambda-e-s\nu$ is an injective linear map on
polynomial coefficients. A fixed left inverse therefore lifts each
old spectral coefficient to a function of the enlarged coefficients.
At $\boldsymbol a=0$ its old-coordinate derivative is unchanged,
and its derivative in the new incoming column is zero. The old
principal moments have the same properties. We need $n$ further
functions with zero old-coordinate
derivatives and independent derivatives in $\boldsymbol a$.

Order the old indices as the $r$ equal-phase indices, the $c$
different-phase indices of the same original group, and the $n$
outside indices:
\begin{equation}\label{eq:canonical-insertion-blocks}
 \mathsf X=
 \begin{pmatrix}e\Id_r&0&\mathsf F\\0&E_C&\mathsf F_C\\\mathsf G&\mathsf G_C&X_{\rm out}\end{pmatrix},
 \qquad
 \mathsf W=\diag\left(\frac1{s-s_h}\right)_{h\ {\rm outside}}.
\end{equation}
Here $e\Id_c-E_C$ is invertible. The growing moment block and its
enlargement are
\[
 \mathsf J=\mathsf F\mathsf W \mathsf G,\qquad
 \mathsf J_+=\begin{pmatrix}\mathsf J&\mathsf F\mathsf W \boldsymbol b\\\boldsymbol a\mathsf W \mathsf G&\boldsymbol a\mathsf W \boldsymbol b\end{pmatrix}.
\]
The principal-moment differences
$\Delta_k=(\tr \mathsf J_+^k-\tr \mathsf J^k)/k$ vanish at $\boldsymbol a=0$, and
\begin{equation}\label{eq:canonical-insertion-moment}
 \left.\mathrm d\Delta_k\right|_{\delta \boldsymbol a}
   =\delta \boldsymbol a\,\mathsf W \mathsf G \mathsf J^{k-2}\mathsf F\mathsf W \boldsymbol b,
 \qquad 2\leq k\leq\min(r,n)+1.
\end{equation}
Newton's identities express the old trace of order $r+1$ in its
first $r$ traces, so every $\Delta_k$ belongs to the enlarged algebra.
The first difference, $\Delta_1$, has derivative
$\delta\boldsymbol a\mathsf W\boldsymbol b$. We do not need it:
the displayed higher differences will already span the moment
directions, and the spectral functions will supply their complement.

First suppose $0<r<n$. Generically $\mathsf J$ is invertible. Set
\begin{equation}\label{eq:canonical-insertion-projection}
 \Pi=\Id_n-\mathsf W \mathsf G \mathsf J^{-1}\mathsf F,\qquad
 \C^n=\ker \mathsf F\oplus\operatorname{im}(\mathsf W \mathsf G).
\end{equation}
The $r$ vectors in \cref{eq:canonical-insertion-moment} span the
second summand if $\mathsf F\mathsf W \boldsymbol b$ is cyclic for $\mathsf J$.
For the other summand use the coefficients of
$R(\nu)=F_+(e+s\nu,\nu)$, where $F_+$ is the enlarged determinant.
They vanish at $\boldsymbol a=0$. With
$\mathsf K(\nu)=(e+s\nu)\Id-\mathsf X-\nu\mathsf D_{\rm old}$,
the Schur complement gives
\[
 R(\nu)=-\boldsymbol a\,
       \operatorname{adj}(\mathsf K(\nu))_{\rm out,out}\boldsymbol b,
 \qquad
 \left.\mathrm dR(\nu)\right|_{\delta \boldsymbol a}
 =-\delta \boldsymbol a\,\det \mathsf K(\nu)\,
            (\mathsf K(\nu)^{-1})_{\rm out,out}\boldsymbol b.
\]
Eliminating the $c$ different-phase indices leaves an outside block
$\nu\mathsf W^{-1}+\mathsf C$, with
\[
 \mathsf C=e\Id_n-X_{\rm out}-\mathsf G_C(e\Id_c-E_C)^{-1}\mathsf F_C,\qquad
 \mathsf R_0=\left.\Pi\mathsf W \mathsf C\right|_{\ker \mathsf F},\quad \boldsymbol c=\Pi\mathsf W \boldsymbol b.
\]
The remaining linear equations give
$(\mathsf K^{-1})_{\rm out,out}\boldsymbol b=(\nu\Id+\mathsf R_0)^{-1}\boldsymbol c$.
To compute its determinant, first eliminate the $c$ different-phase
indices, multiply the outside rows by $\mathsf W$, and use the
decomposition in \cref{eq:canonical-insertion-projection}. The
equal-phase rows pair with $\operatorname{im}(\mathsf W\mathsf G)$
through $\mathsf J$, leaving $\nu\Id+\mathsf R_0$ on $\ker\mathsf F$.
Consequently
\[
 \det\mathsf K(\nu)
 =(-1)^r\det(e\Id_c-E_C)\det(\mathsf W^{-1})\det\mathsf J\,
       \det(\nu\Id+\mathsf R_0).
\]
The prefactor is nonzero and independent of $\nu$.
Hence the coefficient vectors of the new spectral
derivatives span $\ker \mathsf F$ if $\boldsymbol c$ is cyclic for $\mathsf R_0$.
These coefficients are triangular combinations of
$\boldsymbol c,\mathsf R_0\boldsymbol c,\ldots,\mathsf R_0^{n-r-1}\boldsymbol c$.
There are $n-r$ such coefficients, since $\deg R\leq n-r-1$.
When $r=0$, use $\Pi=\Id_n$ and omit the moment differences.

We verify that both cyclicity conditions occur. Choose $\mathsf F$ to select
$r$ outside coordinate directions and $\mathsf G$ so that $\mathsf J$ is diagonal
with distinct nonzero entries. Then $\Pi$ selects the remaining
directions. Condition~\eqref{eq:canonical-insertion-condition} lets
these directions have outside group sizes differing by at most $c+1$.
Pair them across the two groups, leaving at most $c+1$ unpaired
indices. At this test point $\mathsf R_0$ is a direct sum of blocks
\[
 \begin{pmatrix}d_h&a_h\\b_h&d'_h\end{pmatrix}
 \quad\text{and scalar blocks}\quad d_j+\delta_j.
\]
The $d_h,d'_h,d_j$ are fixed by the outside phases; the products
$a_hb_h$ are free. Each paired block can therefore have two distinct
eigenvalues avoiding all eigenvalues already chosen. For all but at most one unpaired
index, use a different-phase index to adjust its diagonal entry
through the term $-\mathsf G_C(e\Id_c-E_C)^{-1}\mathsf F_C$;
these are the independent shifts $\delta_j$. Thus at most one scalar
eigenvalue remains fixed, and the full spectrum can be made simple.
This explains the size hypothesis: even when all unpaired outside
phases coincide, the $c$ available shifts separate all but one of
the at most $c+1$ unpaired directions.
Finally $\boldsymbol b\mapsto(\mathsf F\mathsf W \boldsymbol b,\Pi\mathsf W \boldsymbol b)$ is invertible onto
$\C^r\oplus\ker \mathsf F$, so the two cyclic vectors can be chosen
independently. The required conditions are nonempty Zariski-open
conditions.

If $r\geq n$, choose $\mathsf F,\mathsf G$ so that $\mathsf R=\mathsf W \mathsf G \mathsf F$ has distinct
nonzero eigenvalues. The $n$ vectors in
\cref{eq:canonical-insertion-moment} are now
$\mathsf R^{k-1}\mathsf W \boldsymbol b$, $2\leq k\leq n+1$, and span $\C^n$ for a
cyclic $\mathsf W \boldsymbol b$. No new spectral functions are needed.

In either case the open conditions intersect the old nonzero-Jacobian
locus. The old-coordinate and new-$\boldsymbol a$ derivatives form a block-triangular
Jacobian with nonzero diagonal minors. This proves the insertion step.
\end{proof}

\subsection{Proof of the complete count}

\begin{proof}[Proof of \cref{thm:canonical-independence}]
For $(1,1,1)$ two quadratic spectral coefficients give independent
power combinations. A third contains the sum of the oriented triangle
products, which varies at fixed powers. Their rank is three for every
phase triple.

A positive size triple satisfying the triangle inequalities can be
reduced to $(1,1,1)$ by decreasing a largest entry; the preceding triple
still satisfies those inequalities. Reverse this process, appending
indices last within their equal-phase blocks. At each step
$|m_1-m_2|\leq r+c$, so \cref{lem:canonical-insertion} applies.

If one group exceeds the sum of the others, first build the balanced
triple where its size equals that sum, using initial segments of its
prescribed phase blocks. Append the remaining indices. The growing
group has size $r+c\geq m_1+m_2$, so the lemma applies at every step.
Each insertion adds $n$ canonical pairs and $n$ independent commuting
functions. This proves the full count. Since $H$ commutes with them,
it is locally a function of a complete set and can replace a suitable
member wherever $\mathrm dH\ne0$.
\end{proof}

\subsection{The spectral deficit at repeated phases}

The proof also counts the spectral deficit explicitly. For each
equal-phase block $B$, let $r_B=|B|$, let $o_B$ be the size outside its
original group, and set
\begin{equation}\label{eq:canonical-spectral-deficit}
 k_B=\min(r_B,o_B),\qquad
 q_B=r_Bk_B-\frac{k_B(k_B+1)}2.
\end{equation}
The generic spectral rank is $p-\sum_{B\in\mathcal B}q_B$; the proof
supplies $\sum_{B\in\mathcal B}q_B$ additional independent
principal-moment functions, where $\mathcal B$ is the set of all
equal-phase blocks.
For the lower bound, each insertion increases its block's $q_B$ by
$\min(r,n)$ and leaves the other blocks' $q_B$ unchanged along the
chosen construction. The spectral derivatives supply the complementary
$n-\min(r,n)$ directions. For the upper bound, $\mathsf J_B$ has
$k_B$ distinct nonzero eigenvalues and a semisimple zero eigenvalue of
multiplicity $r_B-k_B$, so its conjugacy orbit has dimension $2q_B$.
The moment-generated tangent space therefore has symplectic rank
$2\sum_{B\in\mathcal B}q_B$. Commuting spectral fields lie in its
symplectic orthogonal, whose maximal isotropic dimension is
$p-\sum_{B\in\mathcal B}q_B$.

For example, with $L=M=2$, $N=1$ and constant phases within each
outer group, there are two size-two equal-phase blocks, each with
$q_B=1$. The spectral rank is six; two additional principal-moment
functions complete the eight-integral count.

